% Bewertungspaket zum Gesamttest «Hydrostatik» — Quelle des PDFs (python3 scripts/build-lp-pdf.py hydrostatik).
% Ganze Punkte je Teilschritt; Werte und Folgefehler-Fälle mit python3 nachgerechnet (05.10.2026).
% Aufbau, Auftrag an die KI und Punktarten (E, W, B, S) wie im Leitprogramm Statik.
\documentclass[11pt]{article}
\usepackage{../lp-druck}
\renewcommand{\lpname}{Hydrostatik}
\renewcommand{\lpteilgebiet}{4.5}
\renewcommand{\lpfuss}{Bewertungspaket}
\newcolumntype{P}{>{\centering\arraybackslash}p{10mm}}
\newenvironment{raster}{\par\smallskip\noindent\tabularx{\linewidth}{@{}>{\raggedright\arraybackslash}p{38mm}P X@{}}\toprule
  \small\textsc{Teilschritt} & \small\textsc{P} & \small\textsc{Musterlösung}\\\midrule}{\bottomrule\endtabularx\par}
\newcommand{\fehler}[1]{\par\smallskip{\small\textcolor{lprot}{\bfseries Typische Fehler:} #1}}
\newcommand{\E}{\,{\footnotesize\bfseries\color{lpgruen}(E)}}
\newcommand{\B}{\,{\footnotesize\bfseries(B)}}
\newcommand{\W}{\,{\footnotesize\bfseries(W)}}
\newcommand{\Sk}{\,{\footnotesize\bfseries(S)}}
\newcommand{\A}{\,{\footnotesize\bfseries\color{lpgruen}(A)}}
\begin{document}
\lpkopf{Bewertungspaket zum Gesamttest}{}

\begin{lpachtung}{Erst lösen, dann öffnen}
Dieses Paket enthält die Musterlösung. Wer es vor dem Lösen liest, misst am Ende nur, wie gut das Abschreiben gelingt.
\end{lpachtung}

\section*{1 · So gehst du vor}
\begin{enumerate}[leftmargin=*,itemsep=2pt]
\item \textbf{Gesamttest lösen} — vollständig, mit Rechenweg, ohne dieses Paket.
\item \textbf{Fotografieren oder scannen}, gut lesbar, eine Seite pro Bild. \textbf{Kein Name} auf den Bildern.
  Handyfotos tragen oft unsichtbar Standort, Zeit und Gerät mit (Metadaten). Lade darum lieber einen Scan oder ein
  Bildschirmfoto des Fotos hoch, oder entferne vorher die Standortangaben.
\item \textbf{Bewerten lassen.} Ob du dafür eine KI nutzt und welche, entscheidest du selbst und in eigener Verantwortung —
  je nachdem, was dir aufgrund deines Alters und deiner persönlichen Situation erlaubt ist (Altersgrenzen und
  Nutzungsbedingungen des Anbieters, allenfalls Einverständnis deiner Eltern). Lade nur deine Lösung und dieses PDF hoch,
  keine persönlichen Daten, und füge den Auftrag aus Abschnitt~2 ein. Ohne KI kannst du dich mit Abschnitt~3 selbst bewerten.
\item \textbf{Rückmeldung prüfen.} Stimmt, was die KI aus deiner Handschrift gelesen hat? Eine KI kann sich irren —
  im Zweifel gilt das Raster in Abschnitt~3.
\item \textbf{Gezielt wiederholen} nach Abschnitt~4.
\end{enumerate}

\needspace{14\baselineskip}
\section*{2 · Auftrag an die KI}
\begin{tcolorbox}[colback=lplinie!25,colframe=lplinie,boxrule=0.5pt,arc=1mm,fontupper=\small]
Du bist Korrektorin oder Korrektor für Physik an einer Schweizer Berufsmaturitätsschule. Im Anhang findest du (1)~das Bewertungspaket zum Gesamttest «Hydrostatik» mit Musterlösung und Punkteraster und (2)~Fotos meiner handschriftlichen Lösung.\par\medskip
Geh so vor:\par
1. Schreib zu jeder Aufgabe G1 bis G6 zuerst ab, was du in meiner Lösung liest — Rechenweg und Ergebnis. Was du nicht sicher lesen kannst, markierst du mit [unsicher]. Nicht raten.\par
2. Vergib die Punkte genau nach dem Raster im Bewertungspaket (Abschnitt 3), ganze Punkte je Zeile. Ziehe einen Fehler nur einmal ab: Rechne ich mit einem falschen Zwischenergebnis richtig weiter, gibt es die Folgepunkte, und derselbe Fehler (etwa dieselbe fehlende Umrechnung) in mehreren Teilaufgaben kostet nur einen Punkt — auch wenn das Folgeergebnis unplausibel ist; eine fehlende Plausibilitätsprobe kostet keinen zusätzlichen Punkt. Die Zeilen «Typische Fehler» sagen, welche Rasterzeile bei diesem Fehler 0 Punkte gibt — sie sind kein zusätzlicher Abzug.\par
3. Andere richtige Lösungswege zählen voll. Gleichwertige Schreibweisen zählen voll: Dezimalpunkt oder Dezimalkomma, andere Einheit oder Vorsilbe ($15.7\;\text{kPa} = 15\,700\;\text{Pa} = 0.157\;\text{bar}$; $4.8\;\text{l} = 4800\;\text{cm}^3 = 0.0048\;\text{m}^3$; $6.4\;\text{g} = 0.0064\;\text{kg}$; $15.3\;\text{t} = 15\,300\;\text{kg}$), Zehnerpotenz oder ausgeschriebene Zahl, sinnvoll gerundete Werte (drei signifikante Stellen). Zeilen mit (E) sind Ergebnispunkte: Sie verlangen das Ergebnis \emph{mit richtiger Einheit}, auch ohne Weg; ohne oder mit falscher Einheit gibt es diesen Punkt nicht. Zeilen mit (A) sind Ablesepunkte: Es zählt der abgelesene Wert mit Einheit, innerhalb der angegebenen Ablesetoleranz. Zeilen mit (W) gibt es nur mit nachvollziehbarem Weg (Formel, dann Werte mit Einheiten); ein Zahlenwert darin braucht ebenfalls die Einheit. Zeilen mit (B) sind Begründungen: Es zählt die fachliche Aussage in eigenen Worten, eine Formel ist nicht nötig; andere richtige Formulierungen zählen voll. Zeilen mit (S) sind Skizzenpunkte: Es zählt, dass die verlangten Kräfte als Pfeile mit Namen und richtiger Richtung eingezeichnet sind; Länge und Schönheit zählen nicht. Werte aus sinnvoll gerundeten Zwischenergebnissen zählen, wenn sie höchstens rund 2\,\% vom Raster abweichen.\par
4. Begründe jeden Abzug in einem Satz und nenne die Fehlerart (zum Beispiel «Fläche nicht in m² umgerechnet»). Schreib die Musterlösung nicht ab — zeig mir, wo mein Fehler liegt.\par
5. Gib zum Schluss eine Tabelle «Aufgabe | Punkte | Begründung», die Gesamtpunktzahl (von 25) und — nach der Selbsteinschätzung in Abschnitt 4 — welches Kapitel des Leitprogramms ich wiederholen soll. Keine Note.\par
6. Fehlt ein Foto oder ist eine Aufgabe unlesbar, sag es, statt sie zu bewerten.
\end{tcolorbox}

\needspace{16\baselineskip}
\section*{3 · Musterlösung und Punkteraster}
\textbf{Allgemein:} Ganze Punkte je Zeile. Ein Fehler wird nur einmal abgezogen: Wer mit einem falschen Zwischenergebnis
richtig weiterrechnet, bekommt die folgenden Zeilen (Folgefehler), und derselbe Fehler in mehreren Teilaufgaben kostet nur einen Punkt — auch wenn das Folgeergebnis unplausibel ist; eine fehlende
Plausibilitätsprobe kostet keinen zusätzlichen Punkt. Gleichwertige Wege und Schreibweisen zählen voll.
In der Spalte P: \textbf{(E)}~= Ergebnispunkt — es zählt das Ergebnis \emph{mit richtiger Einheit}, auch ohne Weg; ohne oder mit
falscher Einheit gibt es ihn nicht. \textbf{(A)}~= Ablesepunkt — der abgelesene Wert mit Einheit, innerhalb der angegebenen Toleranz. \textbf{(W)}~= Wegpunkt — nur mit nachvollziehbarem Weg (Formel, dann Werte mit Einheiten);
steht in einer Weg-Zeile ein Zahlenwert, braucht auch er die Einheit. \textbf{(B)}~= Begründungspunkt — es zählt die fachliche
Aussage in eigenen Worten, eine Formel ist nicht nötig. \textbf{(S)}~= Skizzenpunkt — die verlangten Kräfte als Pfeile mit Namen
und richtiger Richtung. Folgewerte aus gerundeten Zwischenergebnissen zählen, wenn sie höchstens rund 2\,\% abweichen.
«Typische Fehler» nennt die Zeilen mit 0 Punkten und die Punktzahl, die bleibt.

\begin{lpaufgabe}{G1}{Raupenfahrzeug im Diagramm}{4 P}
\begin{raster}
a) Ablesen & 1\A & \(p \approx 50\;\text{kPa}\) bei \(A = 3\;\text{m}^2\) (zählt: \(48\) bis \(52\;\text{kPa}\))\\
b) Kraft, Masse & 1\E & \(F_G = p \cdot A = 50\,000\;\text{Pa} \cdot 3\;\text{m}^2 = 150\,000\;\text{N}\); \(m = \dfrac{F_G}{g} = \dfrac{150\,000\;\text{N}}{9.81\;\text{m/s}^2} \approx 15\,300\;\text{kg} = 15.3\;\text{t}\) (beides nötig; folgerichtig aus dem eigenen Ablesewert)\\
c) Ablesen & 1\A & \(A \approx 6\;\text{m}^2\) bei \(25\;\text{kPa}\) (zählt: \(5.8\) bis \(6.2\;\text{m}^2\)); ebenso gerechnet \(A = \dfrac{F}{p} = \dfrac{150\,000\;\text{N}}{25\,000\;\text{Pa}} = 6\;\text{m}^2\)\\
d) Kurve & 1\B & Die Kraft bleibt gleich, \(p = \dfrac{F}{A}\): doppelte Fläche, halber Druck. Der Druck nimmt darum immer langsamer ab und wird nie null, weil die Kraft immer auf eine endliche Fläche wirkt (umgekehrte Proportionalität, Hyperbel)\\
\end{raster}
\fehler{kPa nicht in Pa umgerechnet (\(F = 150\;\text{N}\), \(m \approx 15.3\;\text{kg}\)): b) 0 \(\to\) 3 von 4\,P. Masse ohne \(g\) gleich der Kraft gesetzt (\(150\;\text{t}\)): b) 0 \(\to\) 3 von 4\,P.}
\end{lpaufgabe}

\begin{lpaufgabe}{G2}{Tauchen im Bergsee, \(p_0 = 840\;\text{hPa}\), \(p = 2.30\;\text{bar}\)}{4 P}
\begin{raster}
Ansatz & 1\W & \(p = p_0 + \rho \cdot g \cdot h\), also \(h = \dfrac{p - p_0}{\rho \cdot g}\), Drücke in Pa\\
a) Tiefe & 1\E & \(h = \dfrac{230\,000\;\text{Pa} - 84\,000\;\text{Pa}}{1000\;\text{kg/m}^3 \cdot 9.81\;\text{m/s}^2} \approx 14.9\;\text{m}\)\\
b) Meer & 1\E & \(h = \dfrac{230\,000\;\text{Pa} - 101\,300\;\text{Pa}}{1025\;\text{kg/m}^3 \cdot 9.81\;\text{m/s}^2} \approx 12.8\;\text{m}\)\\
c) Begründung & 1\B & Im Bergsee drückt weniger Luft auf die Oberfläche, und Süsswasser ist weniger dicht: Für denselben Gesamtdruck braucht es mehr Schweredruck aus einer längeren Wassersäule (beide Gründe oder einer mit klarer Begründung)\\
\end{raster}
\fehler{Luftdruck vergessen (\(h \approx 23.4\;\text{m}\), im Meer \(\approx 22.9\;\text{m}\)): derselbe Fehler in a) und b), Ansatz 0, a) und b) folgerichtig \(\to\) 3 von 4\,P. Im Bergsee mit \(1013\;\text{hPa}\) gerechnet (\(13.1\;\text{m}\)): a) 0 \(\to\) 3 von 4\,P.}
\end{lpaufgabe}

\begin{lpaufgabe}{G3}{Pascals Fassversuch, Fass \(1.6\;\text{m}\), Rohr \(8\;\text{m}\)}{4 P}
\begin{raster}
Ansatz & 1\W & \(p_S = \rho \cdot g \cdot h\), mit \(h\) = Höhe der Wassersäule über dem Boden\\
a) ohne Rohr & 1\E & \(p_S = 1000\;\text{kg/m}^3 \cdot 9.81\;\text{m/s}^2 \cdot 1.6\;\text{m} \approx 15\,700\;\text{Pa} = 15.7\;\text{kPa}\)\\
b) mit Rohr & 1\E & \(h = 1.6\;\text{m} + 8\;\text{m} = 9.6\;\text{m}\): \(p_S = 1000\;\text{kg/m}^3 \cdot 9.81\;\text{m/s}^2 \cdot 9.6\;\text{m} \approx 94\,200\;\text{Pa} = 94.2\;\text{kPa}\), sechsmal so viel\\
c) Begründung & 1\B & Der Schweredruck hängt nur von der Höhe der Wassersäule ab, nicht von der Wassermenge oder der Form des Gefässes (hydrostatisches Paradoxon). Das dünne Rohr macht die Säule hoch, und der grössere Druck wirkt auf die ganze Innenfläche des Fasses\\
\end{raster}
\fehler{Nur das Rohr gerechnet (\(h = 8\;\text{m}\), \(78.5\;\text{kPa}\)): b) 0 \(\to\) 3 von 4\,P. Luftdruck dazugezählt (\(117\;\text{kPa}\) und \(195\;\text{kPa}\)): derselbe Fehler, einmal abgezogen: a) 0, b) folgerichtig \(\to\) 3 von 4\,P. In c) «das Wasser im Rohr ist schwer»: c) 0 — die Menge ist gerade klein.}
\end{lpaufgabe}

\begin{lpaufgabe}{G4}{Wagenheber, \(1200\;\text{kg}\), \(A_2 = 400\;\text{cm}^2\), \(F_1 = 120\;\text{N}\)}{4 P}
\begin{raster}
Ansatz & 1\W & Gleicher Druck an beiden Kolben: \(p = \dfrac{F_2}{A_2} = \dfrac{F_1}{A_1}\), mit \(F_2 = m \cdot g\)\\
a) Druck & 1\E & \(F_2 = 1200\;\text{kg} \cdot 9.81\;\text{m/s}^2 = 11\,772\;\text{N}\); \(p = \dfrac{11\,772\;\text{N}}{0.04\;\text{m}^2} \approx 294\,000\;\text{Pa} \approx 2.94\;\text{bar}\)\\
b) Pumpkolben & 1\E & \(A_1 = \dfrac{F_1}{p} = \dfrac{120\;\text{N}}{294\,300\;\text{Pa}} \approx 4.08 \cdot 10^{-4}\;\text{m}^2 = 4.08\;\text{cm}^2\) (höchstens; gleichwertig \(A_2 \cdot \dfrac{F_1}{F_2}\))\\
c) Öl und Weg & 1\E & \(V = A_2 \cdot s_2 = 400\;\text{cm}^2 \cdot 12\;\text{cm} = 4800\;\text{cm}^3 = 4.8\;\text{l}\); \(s_1 = \dfrac{V}{A_1} = \dfrac{4800\;\text{cm}^3}{4.08\;\text{cm}^2} \approx 1180\;\text{cm} \approx 11.8\;\text{m}\) (beides nötig)\\
\end{raster}
\fehler{Masse statt Gewichtskraft (\(p = 0.3\;\text{bar}\), \(A_1 = 40\;\text{cm}^2\), \(s_1 = 1.2\;\text{m}\)): a) 0, b) und c) folgerichtig \(\to\) 3 von 4\,P. Flächenverhältnis umgekehrt (\(A_1 \approx 39\,200\;\text{cm}^2\)): b) 0, c) folgerichtig \(\to\) 3 von 4\,P.}
\end{lpaufgabe}

\begin{lpaufgabe}{G5}{Ballon, \(12\;\text{l}\), \(8\;\text{g}\)}{5 P}
\begin{raster}
a) Skizze & 1\Sk & Auftrieb senkrecht nach oben, Gewichtskraft senkrecht nach unten, beide mit Namen (der Auftrieb länger gezeichnet schadet nicht)\\
b) Auftrieb & 1\E & \(F_A = \rho_\text{Luft} \cdot V \cdot g = 1.20\;\text{kg/m}^3 \cdot 0.012\;\text{m}^3 \cdot 9.81\;\text{m/s}^2 \approx 0.141\;\text{N}\)\\
c) Gewicht & 1\E & \(F_G = 0.008\;\text{kg} \cdot 9.81\;\text{m/s}^2 \approx 0.0785\;\text{N}\)\\
d) Zusatzlast & 1\E & \(\dfrac{F_A - F_G}{g} = \dfrac{0.0628\;\text{N}}{9.81\;\text{m/s}^2} \approx 0.0064\;\text{kg} = 6.4\;\text{g}\)\\
e) unter Wasser & 1\B & Grösser: Der Ballon verdrängt dasselbe Volumen, aber Wasser ist rund \(800\)-mal dichter als Luft; der Auftrieb hängt an der Dichte des verdrängten Stoffs (eine Zahl ist nicht nötig; nebenbei \(\approx 118\;\text{N}\))\\
\end{raster}
\fehler{Liter nicht in m³ umgerechnet (\(F_A = 141\;\text{N}\)): b) 0; d) folgerichtig zählt, auch wenn das Ergebnis unplausibel ist \(\to\) 4 von 5\,P. Gramm nicht in kg (\(F_G = 78.5\;\text{N}\)): c) 0; d) folgerichtig zählt \(\to\) 4 von 5\,P.}
\end{lpaufgabe}

\begin{lpaufgabe}{G6}{Aräometer, \(30\;\text{g}\), \(2\;\text{cm}^2\)}{4 P}
\begin{raster}
a) Wasser & 1\E & Schwimmen: \(F_A = F_G\), also \(\rho_{Fl} \cdot A \cdot h \cdot g = m \cdot g\) und \(h = \dfrac{m}{\rho_{Fl} \cdot A} = \dfrac{0.03\;\text{kg}}{1000\;\text{kg/m}^3 \cdot 0.0002\;\text{m}^2} = 0.15\;\text{m} = 15\;\text{cm}\)\\
b) Spiritus & 1\E & \(h = \dfrac{0.03\;\text{kg}}{790\;\text{kg/m}^3 \cdot 0.0002\;\text{m}^2} \approx 0.190\;\text{m} = 19.0\;\text{cm}\)\\
c) Skala & 1\B & In einer dichteren Flüssigkeit genügt ein kleineres eingetauchtes Volumen für denselben Auftrieb; das Aräometer taucht weniger tief ein, die Flüssigkeitsoberfläche steht weiter unten am Röhrchen — dort steht darum die grosse Dichte\\
d) Dichte & 1\E & \(\rho_{Fl} = \dfrac{m}{A \cdot h} = \dfrac{0.03\;\text{kg}}{0.0002\;\text{m}^2 \cdot 0.136\;\text{m}} \approx 1100\;\text{kg/m}^3\) (genauer \(1103\))\\
\end{raster}
\fehler{Gramm oder cm² nicht umgerechnet: derselbe Fehler in a), b) und d), einmal abgezogen \(\to\) 3 von 4\,P. Auftrieb mit \(g\), Gewicht ohne \(g\) gerechnet (\(h \approx 1.5\;\text{cm}\)): a) 0, b) und d) folgerichtig \(\to\) 3 von 4\,P.}
\end{lpaufgabe}

\needspace{14\baselineskip}
\section*{4 · Selbsteinschätzung}
\begin{tabularx}{\linewidth}{@{}l X@{}}\toprule
22 – 25 P & Die geprüften Teile sitzen. Wo du Punkte verloren hast: das Kapitel dieser Aufgabe nochmals (Zuordnung unten).\\
17 – 21 P & Den schwächsten Teil nochmals: Simulation und Übungen des Kapitels, in dem du die meisten Punkte verloren hast.\\
11 – 16 P & Zurück zu den Kapiteln aller Aufgaben, in denen du Punkte verloren hast.\\
0 – 10 P & Zurück zu Kapitel 1 und von dort der Reihe nach weiter.\\\midrule
\multicolumn{2}{@{}p{\linewidth}@{}}{\small\textbf{Aufgabe $\to$ Kapitel} (gilt für alle Bereiche): G1 $\to$ Kapitel 1 (K1, K2) · G2 $\to$ Kapitel 2, 3 (K3) · G3 $\to$ Kapitel 2 (K3) · G4 $\to$ Kapitel 4 (K4) · G5 $\to$ Kapitel 5 (K5) · G6 $\to$ Kapitel 6 (K5)}\\\bottomrule
\end{tabularx}
\end{document}
